% Nilpotent truncation fragment.
% Nilpotent truncation route to the classical cotangent shadow.
% Gives a finite classical CE/PV calculation on a precisely specified
% nilpotent truncation of the reduced Hamiltonian Lie algebra \mathfrak h.
% It does not construct the unreduced open BV factorization centre.

\section*{Nilpotent truncation of the cotangent shadow}

The main theorem \ref{thm:main-local} is unconditional for the
cochain-level reduced CE/PV central-operation, principal-part, brane-line factorization,
formal degree-zero Weyl/Moyal coefficient algebra, and scalar-anomaly
statements.  The finite-\(N\) and stable trace comparison through
radial parts remains conditional on the external radial-parts input
named in the degree-zero Moyal theorem.
Part (d) is proved at the reduced principal-part level after de Rham-current contraction;
its unreduced primitive shadow and part (e) on the descendant
(one-antifield) sector are governed by
\ref{prob:formal-factorization-center}, the unreduced Tate-coefficient
cotangent lift. The deciding obstructions are
\ref{lem:tate-casimir-obstruction} (Casimir support fails in the
direct-sum dual) and \ref{lem:finite-window-hamiltonian-obstruction}
(no naive finite-dimensional Lie truncation of the full sector). Two
routes occur in \ref{prob:formal-factorization-center}: a
weighted Tate completion and a nilpotent truncation. The argument below
closes the classical nilpotent-truncation route on a precisely specified
truncation, records what is lost, and records the quantum
Moyal obstruction that prevents the same finite truncation from being
an all-order descendant theorem.

\subsection*{The truncation}

Let $\mathfrak h=(z_1,z_2)\C[[z_1,z_2]]=\mathfrak m$ be the reduced
Hamiltonian Lie algebra of the formal symplectic disk on $\C^2$, with
bracket $\{f,g\}=\partial_{z_1}f\,\partial_{z_2}g-
\partial_{z_2}f\,\partial_{z_1}g$ modulo constants. Let
$\mathfrak m^k=(z_1,z_2)^k\C[[z_1,z_2]]$ for $k\geq 1$ be the
$k$-th power of the maximal ideal. Fix integers $N\geq k\geq 3$ and
define the \emph{nilpotent truncation}
$$\mathfrak h_{\leq N}^{(k)}:=\mathfrak m^k/\mathfrak m^{N+1},$$
with $\mathfrak m$-adic Poisson bracket inherited from the
ambient $\mathfrak h_{\mathrm{poly}}$. The standard choice is
$k=3$:
$$\mathfrak h_{\leq N}:=\mathfrak h_{\leq N}^{(3)}=
  \mathfrak m^3/\mathfrak m^{N+1}.$$
Concretely, $\mathfrak h_{\leq N}$ is the finite-dimensional vector
space spanned by monomials $z_1^az_2^b$ with $3\leq a+b\leq N$.

\begin{rmk}[Why degree three and not degree two]
The naive choice $k=2$ fails to give a nilpotent Lie algebra: the
degree-$2$ piece $H_2=\mathrm{span}\{z_1^2,z_1z_2,z_2^2\}$ is closed
under the Poisson bracket because $\{H_2,H_2\}\subset H_{2+2-2}=H_2$,
and the structure is exactly $\mathfrak{sl}_2$:
$$\{z_1^2,z_2^2\}=4z_1z_2,\quad
  \{z_1^2,z_1z_2\}=2z_1^2,\quad
  \{z_1z_2,z_2^2\}=2z_2^2.$$
Thus $\mathfrak m^2/\mathfrak m^{N+1}$ contains a nonabelian
semisimple subalgebra and is not nilpotent, hence does not satisfy
the lower-central Tate-pronilpotence hypothesis of
\ref{lem:continuous-bar-cobar}. The truncation $k=3$ removes the
$H_2$ direction; the bracket on $\mathfrak m^3$ raises
$\mathfrak m$-adic degree strictly, making the truncation
nilpotent.
\end{rmk}

\begin{lem}[Truncation is a finite-dimensional nilpotent Lie algebra]
\label{lem:trunc-is-nilpotent-lie}
For every $N\geq 3$:
\begin{enumerate}
  \item $\mathfrak h_{\leq N}=\mathfrak m^3/\mathfrak m^{N+1}$ is a
    finite-dimensional Lie algebra of dimension
    $\binom{N+2}{2}-6$ (the number of
    monomials of degree between $3$ and $N$ inclusive).
  \item $\mathfrak h_{\leq N}$ is nilpotent of class at most
    $N-2$ in the lower central series.
  \item $\mathfrak h_{\leq N}^\vee$ is finite-dimensional, hence the
    formal Casimir
    $C_{\mathfrak h_{\leq N}}=\sum_{3\leq d\leq N}\sum_i
     e_{d,i}\otimes e_d^i$
    is a convergent element of
    $\mathfrak h_{\leq N}\otimes\mathfrak h_{\leq N}^\vee$ (a finite
    sum, no Tate completion required).
  \item The linear Hamiltonians $z_1,z_2\in\mathfrak m\setminus
    \mathfrak m^3$, the conformal-symplectic generators
    $H_2=\mathrm{span}\{z_1^2,z_1z_2,z_2^2\}\cong\mathfrak{sl}_2$,
    and the constant Hamiltonian $1\notin\mathfrak m$ are absent
    from $\mathfrak h_{\leq N}$ by construction.
\end{enumerate}
\end{lem}

\begin{proof}
\emph{(i)} The space $\mathfrak m^3$ is a Lie subalgebra of the
polynomial Hamiltonian Lie algebra $\mathfrak h_{\mathrm{poly}}$, and
$\mathfrak m^{N+1}$ is a Lie ideal inside $\mathfrak m^3$: for
$f\in\mathfrak m^p$ and
$g\in\mathfrak m^q$ one has $\{f,g\}\in\mathfrak m^{p+q-2}$, since the
Poisson bracket on $\C[z_1,z_2]$ lowers total polynomial degree by
exactly two, and each partial derivative lowers $\mathfrak m$-adic
degree by exactly one. For $p,q\geq 3$ the bracket lands in
$\mathfrak m^{p+q-2}\subset\mathfrak m^4\subset\mathfrak m^3$, making
$\mathfrak m^3$ a Lie subalgebra; for $p\geq 3$ and $q\geq N+1$
the bracket lands in $\mathfrak m^{p+q-2}\subset\mathfrak m^{N+2}
\subset\mathfrak m^{N+1}$, making $\mathfrak m^{N+1}$ a Lie ideal of
$\mathfrak m^3$. The quotient
$\mathfrak h_{\leq N}=\mathfrak m^3/\mathfrak m^{N+1}$ inherits a
well-defined Lie bracket. Dimension count: monomials $z_1^az_2^b$
with $3\leq a+b\leq N$ number $\binom{N+2}{2}-6$ (the total number
with $0\leq a+b\leq N$ minus the six of degrees $0,1,2$).

\emph{(ii)} The Poisson bracket of two elements of
$\mathfrak m^3$ raises $\mathfrak m$-adic degree by at least $1$:
$\{H_p,H_q\}\subset H_{p+q-2}\subset\mathfrak m^{p+q-2}$, and for
$p,q\geq 3$, $p+q-2\geq 4>p$ (since $q\geq 3$). Thus the descending
Lie series
$$\mathfrak h_{\leq N}^{(\geq m)}:=
  \underbrace{[\mathfrak h_{\leq N},[\mathfrak h_{\leq N},\ldots]]}_{m\text{ brackets}}
  \subset\bigl(\mathfrak m^{m+3}/\mathfrak m^{N+1}\bigr)$$
satisfies $\mathfrak h_{\leq N}^{(\geq m)}=0$ for
$m+3\geq N+1$, i.e.\ for $m\geq N-2$. Hence $\mathfrak h_{\leq N}$
is nilpotent of class at most $N-2$.

\emph{(iii)} Immediate from finite-dimensionality.

\emph{(iv)} Immediate from the $\mathfrak m^3$-adic cutoff and the
exclusion $a+b\leq 2$.
\end{proof}

\begin{rmk}[Why the maximal-ideal-adic obstruction does not apply]
\label{rmk:why-no-finite-window-obstruction}
\ref{lem:finite-window-hamiltonian-obstruction} shows that
$\mathfrak h/\mathfrak m^{N+1}$ is \emph{not} a Lie quotient of the
full Hamiltonian algebra $\mathfrak h_{\mathrm{poly}}$: the linear
Hamiltonian $z_1$ acting on $z_2^{N+1}\in\mathfrak m^{N+1}$ produces
$\{z_1,z_2^{N+1}\}=(N+1)z_2^N\notin\mathfrak m^{N+1}$. The
nilpotent truncation $\mathfrak h_{\leq N}=\mathfrak m^3/
\mathfrak m^{N+1}$ \emph{kills the linear Hamiltonian} and the
conformal $\mathfrak{sl}_2$ direction from the
algebra. The obstruction lemma is then irrelevant: no element of
$\mathfrak h_{\leq N}$ has $\mathfrak m$-adic degree below $3$, so the
displayed counterexample $\{z_1,\,z_2^{N+1}\}$ never occurs in
$\mathfrak h_{\leq N}$.
\end{rmk}

\begin{rmk}[Why the Tate Casimir obstruction does not apply]
\label{rmk:why-no-tate-casimir}
The product/direct-sum support obstruction of
\ref{lem:tate-casimir-obstruction} is a feature of the unbounded
expansion $\mathfrak h=\prod_{d\geq 1}H_d$: the Casimir
$C_{\mathfrak h}=\sum_d\sum_i e_{d,i}\otimes e_d^i$ has nonzero
component in every $H_d^\vee$, and $\bigoplus_d H_d^\vee$ admits no
infinite-support element. On the truncation $\mathfrak h_{\leq N}$
the Casimir is the finite sum
$C_{\mathfrak h_{\leq N}}=\sum_{d=3}^N\sum_i e_{d,i}\otimes e_d^i$
in the finite-dimensional vector space
$\mathfrak h_{\leq N}\otimes\mathfrak h_{\leq N}^\vee$. There is
no convergence question.
\end{rmk}

\subsection*{The classical truncated cotangent CE/PV shadow}

\begin{thm}[Classical truncated cotangent CE/PV shadow]\label{thm:phi-trunc-classical}
Set
$\mathfrak g_{\leq N}=\mathfrak h_{\leq N}\ltimes
\mathfrak h_{\leq N}^\vee[1]$. For every interval $I\subset\R$ let
$\mathfrak g_{\leq N,I}=
\Omega^\bullet_c(I)\,\widehat\otimes\,\mathfrak h_{\leq N}\,\ltimes\,
\mathcal D'_c(I)\,\widehat\otimes\,\mathfrak h_{\leq N}^\vee[1]$.
Then, after the Poisson, residue, and interval-smearing conventions fixed
above, there is a filtered cochain isomorphism of completed
dg graded-commutative $P_0$-algebras
$$\Phi_{\leq N}:
  C^\bullet_{\mathrm{CE},\mathrm{cont}}(\mathfrak g_{\leq N,I})
  \xrightarrow{\ \sim\ }
  \mathrm{PV}\bigl(\widehat{\mathbf S}(\mathfrak h_{\leq N,I})\bigr)$$
on every interval \(I\subset\R\). This is the finite nilpotent
CE/PV cotangent shadow of the reduced Hamiltonian model; it is not an
unreduced open BV factorization-centre morphism. The isomorphism is determined on generators by
$\Phi_{\leq N}(c^I)=\theta^I$, $\Phi_{\leq N}(u_I)=O_I$, with the
same Schouten/Lichnerowicz convention as in
\ref{thm:reduced-ce-pv-central-operation}. It is compatible with extension
by zero and with disjoint-interval products in this reduced finite
model, and on the Hamiltonian generators it recovers the stalkwise
operation \(f\mapsto L_{B_f}\). The cotangent generators are finite
principal-part labels in the CE/PV shadow; their unreduced boundary
centrality homotopies remain part of
Problem~\ref{prob:formal-factorization-center}.
\end{thm}

\begin{proof}
The Lie algebra $\mathfrak g_{\leq N}$ is finite-dimensional
($\dim\mathfrak g_{\leq N}=2\dim\mathfrak h_{\leq N}=
2\binom{N+2}{2}-12$), hence pro-finite-dimensional in the trivial
constant filtration. It is lower-central Tate-pronilpotent in the sense of
\ref{lem:continuous-bar-cobar}: by
\ref{lem:trunc-is-nilpotent-lie}(ii) the descending Lie powers of
$\mathfrak h_{\leq N}$ vanish after at most $N-2$ brackets.  The
cotangent summand is abelian, but this alone would not imply nilpotence
of the semidirect product.  The needed point is the degree-lowering
coadjoint action: since \([H_p,H_q]\subset H_{p+q-2}\), an element
\(f\in H_p\), \(p\geq3\), sends \(H_d^\vee\) to
\(H_{d-p+2}^\vee\), hence lowers dual degree by at least one.  Repeated
action therefore kills every finite quotient
\(\mathfrak h_{\leq N}^\vee\).  Together with nilpotence of
\(\mathfrak h_{\leq N}\), the descending Lie powers of
\(\mathfrak g_{\leq N}=\mathfrak h_{\leq N}\ltimes
\mathfrak h_{\leq N}^\vee[1]\) also terminate. Hence
$\bigcap_m\mathfrak g_{\leq N}^{(\geq m)}=0$ in any Tate
filtration, and the lower-central Tate-pronilpotence hypothesis of
\ref{lem:continuous-bar-cobar} is satisfied trivially.

Apply \ref{thm:reduced-ce-pv-central-operation} (the cochain-level reduced
Hamiltonian CE/PV central-operation theorem) to $\mathfrak h_{\leq N}$.
The conclusion is a canonical isomorphism of completed dg
graded-commutative $P_0$-algebras
$$\Phi_{\leq N}^{\mathrm{stalk}}:
  C^\bullet_{\mathrm{CE},\mathrm{cont}}(\mathfrak g_{\leq N})
  \xrightarrow{\ \sim\ }
  \mathrm{PV}\bigl(\widehat{\mathbf S}(\mathfrak h_{\leq N})\bigr),$$
with the explicit generator-level data of
\ref{thm:reduced-ce-pv-central-operation}.

To pass from the stalk identification to the factorization-algebra
identification on intervals $I\subset\R$, the Hamiltonian-sector
factorization upgrade of \ref{thm:hamiltonian-current-factorization}
applies verbatim with $\mathfrak h$ replaced by
$\mathfrak h_{\leq N}$: the locality input
$\{(\phi_1)^i{}_j(t),(\phi_2)^k{}_l(s)\}=\delta^i_l\delta^k_j\delta(t-s)$
is unchanged, the smearing
  $\mathfrak h_{\leq N,I}=\Omega^\bullet_c(I)\,\widehat\otimes\,
  \mathfrak h_{\leq N}$ is canonical because $\mathfrak h_{\leq N}$
is finite-dimensional (no Tate completion), and the strict $P_0$
bracket follows from the delta-function locality.

In this finite truncated classical sector, the cotangent factor has no
Tate-Casimir convergence problem. The principal-part sub-factorization algebra
$$\widetilde{\mathrm{Obs}}^{\mathrm{cl,pp}}_{\partial,\leq N}(I)
  =\mathrm{Obs}^{\mathrm{cl}}_{\mathrm{open},\leq N}(I)
   \,\widehat\otimes\,
   \widehat{\mathbf S}\bigl(\mathfrak h_{\leq N,I}^\vee[1]\bigr)$$
    admits a chain-level primitive projection
    $\pi_{\mathrm{prim}}^{\leq N}$ because the Tate Casimir
    $C_{\mathfrak h_{\leq N}}=\sum_{d=3}^N\sum_i e_{d,i}\otimes e_d^i$
is a finite sum, hence a literal element of the (uncompleted) tensor
product $\mathfrak h_{\leq N}\otimes\mathfrak h_{\leq N}^\vee$. The
$Q$-compatibility, factorization compatibility, and $P_0$-bracket
homotopies are then formal, since at no step does one need an
infinite-support tensor in the direct-sum dual.

The bracket compatibility and graded Leibniz extension to the completed
finite CE/PV algebra is the same argument as in
\ref{thm:reduced-ce-pv-central-operation}, with the additional
simplification that all completions are trivial because
$\mathfrak h_{\leq N}$ is finite-dimensional. Compatibility with
extension by zero, disjoint-interval factorization, and the
stalkwise limit follow as in
\ref{thm:hamiltonian-current-factorization}.
\end{proof}

\subsection*{Quantum obstruction for the finite truncation}

\begin{thm}[Quantum obstruction for the nilpotent truncation]
\label{thm:phi-hbar-trunc}
  The classical finite truncation
  \ref{thm:phi-trunc-classical} is unconditional.  The same finite
  Lie algebra
  $\mathfrak h_{\leq N}=\mathfrak m^3/\mathfrak m^{N+1}$ does
  \emph{not} carry an all-order Moyal deformation for $N\geq4$:
  it is not closed under the reduced Moyal bracket. Consequently the
  nilpotent truncation
  supplies no all-order quantum descendant theorem on the strict
  finite truncation unless an additional Moyal-flat replacement or a
  projected bracket with Jacobi and cochain compatibility is supplied.

  What remains true at quantum level is the restriction of the
  full-sector degree-zero map
  \ref{thm:phi-hbar-all-order} to the finite set of generators of
  $\mathfrak h_{\leq N}$, viewed inside the full
  $\mathfrak h_\hbar$.  This is a restriction of a proved map on the
  full Hamiltonian source, not a closed CE theorem for a finite
  quantum Lie subalgebra.
\end{thm}

\begin{proof}
  The Moyal bracket is
  $$\{f,g\}_\hbar=\sum_{j\geq 0}
    \frac{\hbar^{2j}}{(2j+1)!\,2^{2j}}P^{2j+1}(f,g),\qquad
    P=\partial_{z_1}\otimes\partial_{z_2}
      -\partial_{z_2}\otimes\partial_{z_1}.$$
  If $f\in\mathfrak m^p$ and $g\in\mathfrak m^q$, the
  order-$\hbar^{2j}$ summand lies in
  $\mathfrak m^{p+q-2(2j+1)}$ when nonzero.  The Poisson term
  $j=0$ preserves $\mathfrak m^3$ as a Lie subalgebra, but the
  higher Moyal terms need not.  For $N\geq4$,
  $$P^3(z_1^4,z_2^4)=24z_1\cdot24z_2=576z_1z_2,$$
  so
  $$\{z_1^4,z_2^4\}_\hbar
    =16z_1^3z_2^3+24\hbar^2 z_1z_2+\cdots.$$
  The term $24\hbar^2z_1z_2$ has $\mathfrak m$-adic degree $2$.
  It is neither in $\mathfrak m^3$ nor a constant killed by the
  reduction modulo $\C\cdot1$.  Thus
  $\mathfrak m^3/\mathfrak m^{N+1}$ is not a Moyal Lie subalgebra
  for $N\geq4$.

  Therefore the CE differential for a hypothetical finite quantum
  source
  $\mathfrak h_{\leq N,\hbar}\ltimes
  \mathfrak h_{\leq N,\hbar}^\vee[1]$ is not defined by restriction
  of the Moyal bracket.  Any projected bracket would require a new
  Jacobi proof; the projected-degree warning of
  \ref{lem:finite-window-hamiltonian-obstruction} shows that this is
  a real condition, not a formality.  The all-order theorem
  \ref{thm:phi-hbar-all-order} remains valid on the full
  Hamiltonian algebra $\mathfrak h_\hbar$, and one may evaluate its
  map on generators of degrees $3,\ldots,N$.  That operation does not
  make those generators a finite quantum CE source.  The descendant
  cotangent sector at $\hbar=0$ is exactly
  \ref{thm:phi-trunc-classical}; at $\hbar\neq0$ it remains a named
  Moyal-flat truncation problem.
\end{proof}

\subsection*{What is lost}

\begin{prop}[Lost sectors]\label{prop:trunc-lost-sectors}
The truncation $\mathfrak h\rightsquigarrow\mathfrak h_{\leq N}=
\mathfrak m^3/\mathfrak m^{N+1}$ removes the following five
data from the local theorem \ref{thm:main-local}:
\begin{enumerate}
  \item\label{lost:translations}
    The linear Hamiltonians $z_1,z_2\in\mathfrak m\setminus\mathfrak m^2$,
    which generate translations in the symplectic plane and the
    center-of-mass mode of the brane.
  \item\label{lost:sl2}
    The conformal-symplectic $\mathfrak{sl}_2$ generators
    $H_2=\mathrm{span}\{z_1^2,z_1z_2,z_2^2\}$, which generate
    rotations and dilations of the symplectic plane. This is the
    semisimple sector that prevents $\mathfrak m^3/\mathfrak m^{N+1}$
    from being nilpotent.
  \item\label{lost:constants}
    The constant Hamiltonian $1\notin\mathfrak m$, on which the
    $\hbar$-normalization of the Moyal product depends. The
    central element $\C\cdot 1\subset\mathfrak h_{\mathrm{poly}}$
    is absent from $\mathfrak m^3$ a fortiori, so the
    $\hbar$-normalization in the truncated theorem is conventional.
  \item\label{lost:u1-anomaly}
    The U(1) center anomaly cohomology class
    $[\bar c]\in H^2_{\mathrm{Lie}}(\bar A,\C)$ of
    \ref{thm:main-local}(f), since the cocycle
    $$\omega(z_1^kz_2^l,z_1^mz_2^n)=(kn-lm)\,
       \mathbf 1[(k+m,l+n)=(1,1)]$$
    is concentrated on multidegree $(1,1)$, with both arguments of
    total $\mathfrak m$-adic degree $1$, hence the cocycle pairs
    only elements outside $\mathfrak h_{\leq N}=\mathfrak m^3/
    \mathfrak m^{N+1}$. There is no central anomaly to anchor on
    the truncation.
  \item\label{lost:colim}
    The stable connected limit $N\to\infty$ on the full
    $\mathfrak h$. The inverse system
    $\{\mathfrak h_{\leq N}\}_{N\geq 3}$ has structure maps the
    natural surjections
    $\mathfrak h_{\leq N+1}\twoheadrightarrow\mathfrak h_{\leq N}$
    induced by the $\mathfrak m$-adic projection
    $\mathfrak m^{N+1}\subset\mathfrak m^{N+2}$. The inverse limit
    is $\varprojlim_N\mathfrak h_{\leq N}=\mathfrak m^3$, the
    $\mathfrak m^3$-truncation of the Hamiltonian Lie algebra, not
    the full $\mathfrak h=\mathfrak m$.
\end{enumerate}
What remains classically: the third-order and higher canonical
transformations, which form a nilpotent Lie subalgebra
$\mathfrak m^3\subset\mathfrak h$, a Lie ideal in $\mathfrak m^2$,
admitting the classical finite cotangent CE/PV shadow unconditionally.
The unreduced open BV factorization-centre morphism is not constructed
by this truncation. The all-order Moyal descendant lift is excluded by
\ref{thm:phi-hbar-trunc} unless extra quantum truncation data are
added.
\end{prop}

\begin{proof}
(i)--(iv) are immediate from the definitions and from the
$\mathfrak m$-adic concentration of $\omega$. (v) is the explicit
inverse limit calculation: a compatible family
$\{f_N\in\mathfrak m^3/\mathfrak m^{N+1}\}_N$ assembles to a unique
formal series in $\mathfrak m^3\subset\C[[z_1,z_2]]$, and the
$\mathfrak m$-adic Lie structure is preserved at every level.
\end{proof}

\subsection*{Cubic inverse-limit survivor}

\begin{prop}[Cubic coordinate-window survivor and variance boundary]
\label{prop:colim-recovery}
The finite Lie quotients
\[
  \mathfrak h_{\leq N}=\mathfrak m^3/\mathfrak m^{N+1}
\]
recover the complete cubic Lie algebra
\(\varprojlim_N\mathfrak h_{\leq N}=\mathfrak m^3\).  The finite
cotangent maps \(\Phi_{\leq N}\) of
\ref{thm:phi-trunc-classical} form a compatible family of coordinate
tests.  They do not, however, form an inverse system of dg CE cochain
projections.  A Lie quotient
\(\mathfrak h_{\leq M}\twoheadrightarrow\mathfrak h_{\leq N}\) induces CE
cochain maps by pullback in the opposite direction.  Consequently
\(\Phi_{\mathfrak m^3}\) denotes the coordinate-window family
\(\{\Phi_{\leq N}\}_N\) together with the chain-side inverse limit and an
exact continuous-dualization datum; it is not an unconditional formula
\(\varprojlim_N\Phi_{\leq N}\) on CE cochains.

This is the maximal survivor of the nilpotent tower only in that precise
sense: the Lie-side tower recovers \(\mathfrak m^3\), while a completed
cochain \(P_0\)-statement requires the additional variance and
bracket-admissibility data named above.
\end{prop}

\begin{proof}
The Lie-side inverse limit is the usual \(\mathfrak m\)-adic completion
of \(\mathfrak m^3\).  The variance warning is forced already at
\(M=4\), \(N=3\).  In \(\mathfrak h_{\leq4}\),
\[
  \{z_1^3,z_2^3\}=9z_1^2z_2^2,
  \qquad
  d_{\mathrm{CE}}c^{2,2}\supset -9c^{3,0}c^{0,3}.
  \]
The coordinate projection to the \(N=3\) cochain window kills
\(c^{2,2}\) but not \(c^{3,0}c^{0,3}\), so
\[
  p(d_{\mathrm{CE}}c^{2,2})=-9c^{3,0}c^{0,3}
  \neq0=d_{\mathrm{CE}}p(c^{2,2}).
  \]
Thus projection of CE cochains is not a dg map.  The correct finite
functoriality is contravariant pullback along Lie quotients, while a
completed cochain theorem must be obtained either from a chain-side
inverse limit followed by exact continuous dualization or from a
separately supplied topology whose cochain transition maps are dg maps.

The Lie algebra $\mathfrak m^3$ is lower-central Tate-pronilpotent in the
$\mathfrak m$-adic Tate filtration: its descending Lie powers satisfy
$(\mathfrak m^3)^{(\geq m)}\subset\mathfrak m^{m+3}$ because each
iterated bracket raises the $\mathfrak m$-adic degree by at least
$1$ when both arguments lie in $\mathfrak m^3$ (since
$\{H_p,H_q\}\subset\mathfrak m^{p+q-2}$ and $p,q\geq 3$ implies
$p+q-2\geq p+1$). Hence $\bigcap_m(\mathfrak m^3)^{(\geq m)}=0$,
	and the lower-central Tate-pronilpotence hypothesis of
\ref{lem:continuous-bar-cobar} holds. The Tate Casimir
$C_{\mathfrak m^3}=\sum_{d\geq 3}\sum_i e_{d,i}\otimes e_d^i$
is governed by the same product/direct-sum convergence question
as in \ref{lem:tate-casimir-obstruction}, but on the smaller
	direct-sum dual $\bigoplus_{d\geq 3}H_d^\vee\subset
\bigoplus_{d\geq 1}H_d^\vee$.  The finite maps
\(\Phi_{\leq N}\) at every finite \(N\) use only the finite Casimir
\(\sum_{3\leq d\leq N}\sum_i e_{d,i}\otimes e_d^i\).  Passing from these
finite identities to a completed cochain identity is exactly the
continuous-dualization and variance datum isolated in the statement, not
a formal inverse-limit consequence.
\end{proof}

\begin{defn}[Low-degree obstruction datum for the nilpotent tower]
\label{defn:nilpotent-low-degree-obstruction-terms}
  For \(N\geq3\) let
  \[
    \pi_N\colon\mathfrak m^3\twoheadrightarrow
    \mathfrak h_{\leq N}=\mathfrak m^3/\mathfrak m^{N+1}
  \]
  be the nilpotent window projection.  The two obstruction terms for
  adjoining the removed low-degree sector are:
  \[
    \mathfrak o^{\mathrm{lin}}_N(x;\tau)
      =\{x,\tau\}\bmod \mathfrak m^{N+1},
      \qquad
      x\in H_1,\ \tau\in \mathfrak m^{N+1}/\mathfrak m^{N+2},
  \]
  and
  \[
    \mathfrak o^{\mathrm{ss}}\colon\Lambda^2H_2\to H_2,\qquad
    \mathfrak o^{\mathrm{ss}}(x,y)=\{x,y\}.
  \]
  Concretely
  \[
    \mathfrak o^{\mathrm{lin}}_N(z_1;z_2^{N+1})
      =(N+1)z_2^N\neq0
      \quad\text{in }\mathfrak m^N/\mathfrak m^{N+1},
  \]
  and
  \[
    \{z_1^2,z_2^2\}=4z_1z_2,\qquad
    \{z_1^2,z_1z_2\}=2z_1^2,\qquad
    \{z_1z_2,z_2^2\}=2z_2^2 .
  \]
  The compatible Casimir missing from the nilpotent tower is
  \[
    \mathfrak o^{\mathrm{Cas}}_{\leq2}
      =\sum_i e_{1,i}\otimes e_1^i
       +\sum_i e_{2,i}\otimes e_2^i .
  \]
\end{defn}

\begin{lem}[Obstruction-datum criterion for enlarging the tower]
\label{lem:nilpotent-obstructions-exact}
  A finite \(\mathfrak m\)-adic tower extending
  \(\{\mathfrak h_{\leq N}\}_{N\geq3}\) to a Lie tower containing
  \(H_1\oplus H_2\) and using the same quotient maps exists only if
  \(\mathfrak o^{\mathrm{lin}}_N=0\) for all \(N\) and
  \(\mathfrak o^{\mathrm{ss}}\) is nilpotent in the lower-central
  filtration.  A cotangent CE/PV extension of the inverse limit to the
  full \(\mathfrak h=\mathfrak m\) also requires a coefficient category
  in which the compatible tensor
  \(\mathfrak o^{\mathrm{Cas}}_{\leq2}
    +\sum_{d\geq3}\sum_i e_{d,i}\otimes e_d^i\)
  converges.
\end{lem}

\begin{proof}
  If \(H_1\) is included and the window \(N\) is still the quotient by
  \(\mathfrak m^{N+1}\), then the tail
  \(\mathfrak m^{N+1}\) must be a Lie ideal for the enlarged source.
  The displayed value
  \(\{z_1,z_2^{N+1}\}=(N+1)z_2^N\notin\mathfrak m^{N+1}\)
  shows that this ideal condition is exactly the vanishing of
  \(\mathfrak o^{\mathrm{lin}}_N\), and it fails.

  If \(H_2\) is included, then its bracket remains inside \(H_2\) and
  equals the nonabelian \(\mathfrak{sl}_2\) bracket displayed in
  Definition~\ref{defn:nilpotent-low-degree-obstruction-terms}.  The
  image of \(\mathfrak o^{\mathrm{ss}}\) is all of \(H_2\), so the
  descending Lie powers never leave \(H_2\).  This is exactly the
  obstruction to nilpotence and to the conilpotent bar-cobar hypothesis.

  The finite cotangent windows for
  \(\mathfrak h_{\leq N}\) use the finite Casimir
  \(\sum_{3\leq d\leq N}\sum_i e_{d,i}\otimes e_d^i\).  Extending to
  \(\mathfrak h\) adds the missing \(d=1,2\) components and then the
  whole compatible family over all \(d\geq1\).  Thus the full
  cotangent lift requires precisely the convergence of the displayed
  full Casimir family.
\end{proof}

\begin{thm}[Low-degree obstruction datum for full Hamiltonian recovery]
\label{rmk:colim-does-not-recover-full}
The coordinate-window family
  $\Phi_{\mathfrak m^3}$ of \ref{prop:colim-recovery} is the
  classical cotangent lift datum on $\mathfrak m^3$, not on
$\mathfrak h=\mathfrak m$. To extend to all of $\mathfrak h$ one must
include the low-degree sectors
$\mathfrak m/\mathfrak m^3=H_1\oplus H_2=\C\cdot z_1\oplus
\C\cdot z_2\oplus\mathrm{span}\{z_1^2,z_1z_2,z_2^2\}$, which by
\ref{lem:finite-window-hamiltonian-obstruction}(iii) and
\ref{lem:nilpotent-obstructions-exact} is not stable
under any $\mathfrak m$-adic cutoff: the bracket
$\{z_1,z_2^{N+1}\}=(N+1)z_2^N$ is the explicit obstruction.
Equivalently, the obstruction datum is
\[
  \bigl(\mathfrak o^{\mathrm{lin}}_N,\,
    \mathfrak o^{\mathrm{ss}},\,
    C_{\mathfrak h}\bigr),
\]
where
\[
  \mathfrak o^{\mathrm{lin}}_N(z_1;z_2^{N+1})
    =(N+1)z_2^N\neq0,\qquad
  \operatorname{im}\mathfrak o^{\mathrm{ss}}=H_2\cong\mathfrak{sl}_2,
\]
and \(C_{\mathfrak h}\) is the full infinite Casimir family below.
Equivalently, the Tate Casimir
$C_{\mathfrak h}=\sum_{d\geq 1}\sum_i e_{d,i}\otimes e_d^i$ has
low-degree components
$\sum_i e_{1,i}\otimes e_1^i\in H_1\otimes H_1^\vee$ and
$\sum_i e_{2,i}\otimes e_2^i\in H_2\otimes H_2^\vee$
that the truncation removes; the degree-$1$ component is the source
of the linear-Hamiltonian translation, and the degree-$2$ component
is the source of the conformal-symplectic $\mathfrak{sl}_2$ direction
on which the truncation fails to be nilpotent. Recovering the full
$\mathfrak h$ requires a coefficient category in which
$C_{\mathfrak h}$ converges \emph{including} both low-degree
components, which is exactly the Casimir support obstruction of
\ref{lem:tate-casimir-obstruction}. The truncation route does not
solve that obstruction; it removes the sector on which the
obstruction is felt.  A full recovery theorem must therefore supply
both a coefficient category in which \(C_{\mathfrak h}\) converges and a
relative or reductive model for the \(H_2\)-sector.
\end{thm}

\begin{proof}
  The tower used in \ref{prop:colim-recovery} is
  \(\mathfrak h_{\leq N}=\mathfrak m^3/\mathfrak m^{N+1}\).
  Since the transition maps are the usual quotient maps,
  \[
    \varprojlim_N \mathfrak h_{\leq N}
      =\varprojlim_N \mathfrak m^3/\mathfrak m^{N+1}
      =\mathfrak m^3
  \]
  as a complete filtered Lie algebra.  The homogeneous pieces
  \(H_1\) and \(H_2\) have zero image in every term of this tower;
  hence no element of the inverse limit can recover them.

  The failure is not an artefact of the chosen notation.  If one tries
  to include \(H_1\) while keeping the same \(\mathfrak m\)-adic
  finite windows, the tails are not Lie ideals: for every \(N\),
  \[
    z_2^{N+1}\in\mathfrak m^{N+1},\qquad
    \{z_1,z_2^{N+1}\}=(N+1)z_2^N\notin\mathfrak m^{N+1}.
  \]
  Therefore the quotient by the tail does not carry the full
  Hamiltonian bracket.  If one includes \(H_2\), the degree-two
  Hamiltonians contain the \(\mathfrak{sl}_2\) bracket
  \([H_2,H_2]=H_2\), so the finite quotient is not nilpotent.  Thus a
  finite nilpotent \(\mathfrak m\)-adic tower cannot contain the
  low-degree sector \(H_1\oplus H_2\).

  Finally, the full Casimir \(C_{\mathfrak h}\) has nonzero components
  in \(H_1\otimes H_1^\vee\), \(H_2\otimes H_2^\vee\), and in every
  higher \(H_d\otimes H_d^\vee\).  The nilpotent tower sees only the
  \(d\geq3\) part.  Hence the inverse-limit cotangent lift on
  \(\mathfrak m^3\) cannot be promoted to the full Hamiltonian sector
  without an additional coefficient category in which the full Casimir
  converges.
\end{proof}

\begin{cor}[Full maximal-ideal universal obstruction class]
\label{cor:full-maximal-ideal-universal-obstruction-class}
  For the full reduced formal Hamiltonian algebra
  \(\mathfrak h=\mathfrak m=(z_1,z_2)\C[[z_1,z_2]]\), the direct
  formal-local CE/PV recognition theorem remains valid in the
  finite-coordinate product completion.  The universal filtered-cobar
  Koszul recognition criterion of
  Theorem~\ref{thm:universal-filtered-cobar-recognition}, however, has
  a nonzero obstruction class in
  \(\mathcal O_{\mathrm{Kos}}\) of
  Definition~\ref{defn:universal-koszul-obstruction-complex}.  Its
  visible components are
  \[
    \mathfrak o^{\mathrm{lin}}_N(z_1;z_2^{N+1})
      =(N+1)z_2^N\neq0,
    \qquad
    \operatorname{im}\mathfrak o^{\mathrm{ss}}=H_2\cong\mathfrak{sl}_2,
  \]
  \[
    [C_{\mathfrak h}]
      =
      \left[
        \sum_{d\geq1}\sum_i e_{d,i}\otimes e_d^i
      \right]
      \in Q_{\mathrm{Cas}},
    \qquad
    H_2\subset\bigcap_{r\geq1}\mathfrak h^{(r)} .
  \]
  Therefore a universal full-\(\mathfrak h\)
  bar-cobar/enveloping theorem requires,
  before any open-side \(A_\infty\) comparison, all of the following:
  a coefficient topology representing \(C_{\mathfrak h}\), continuity of
  the finite-window bracket and coadjoint transition maps, vanishing
  \(\varprojlim^1\) for the completed CE/PV and cobar towers, a
  relative or reductive treatment of the
  \(H_2\)-sector, and a twisting cochain whose associated-graded cobar
  map has acyclic cone.
\end{cor}

\begin{proof}
  The formal-local CE/PV part is
  Theorem~\ref{thm:formal-disk-completed-ce-pv}, equivalently the
  formal-disk instance of
  Theorem~\ref{thm:universal-formal-local-ce-pv-recognition}: the
  bracket and coadjoint formulas are continuous for finite coordinate
  projections, so the generator substitution
  \(c\mapsto\theta,\ u\mapsto O\) is a completed dg commutative cochain
  isomorphism.  The \(P_0\)-algebra reading is finite-level or
  bracket-admissible data; it is not supplied by an unqualified inverse
  limit of ambient coordinate products.

  The filtered-cobar recognition theorem asks for more.  The linear
  obstruction is the same computation as in
  Definition~\ref{defn:nilpotent-low-degree-obstruction-terms}:
  adjoining \(H_1\) to the \(\mathfrak m\)-adic nilpotent windows makes
  the tail non-invariant because
  \(\{z_1,z_2^{N+1}\}=(N+1)z_2^N\notin\mathfrak m^{N+1}\).  The
  semisimple obstruction is the displayed \(\mathfrak{sl}_2\) bracket
  on \(H_2\), hence \(H_2\subset\bigcap_r\mathfrak h^{(r)}\).  This
  gives a nonzero lower-central component in
  \(\mathcal O_{\mathrm{conilp}}\).

  The Casimir component is nonzero in the strict product/direct-sum
  topology by the support argument of
  Theorem~\ref{thm:lane-strict-casimir-obstruction}: the compatible
  finite-window Casimirs have nonzero identity component in every
  \(H_d\otimes H_d^\vee\), while a strict tensor with a direct-sum dual
  has only finitely many dual-degree components.  Finally, even after
  changing topology and handling \(H_2\) relatively, the universal
  criterion still demands the cobar cone and the Milnor
  \(\varprojlim^1\) terms in
  Definition~\ref{defn:universal-koszul-obstruction-complex} to vanish.
  These are precisely the listed remaining requirements.
\end{proof}

\begin{prob}[Low-degree-sector lift from the cubic ideal to the maximal ideal]
\label{prob:linear-sector-lift}
Extend the coordinate-window cotangent lift datum $\Phi_{\mathfrak m^3}$ of
\ref{prop:colim-recovery} from $\mathfrak m^3$ to all of
$\mathfrak h=\mathfrak m$. Equivalently, supply a Tate or
weighted-coefficient extension in which the degree-$1$ and
degree-$2$ Casimir components converge as kernels for the BV
propagator, and in which the conformal-symplectic
$\mathfrak{sl}_2$ direction $H_2$ is included by a relative or reductive
replacement for the lower-central Tate-pronilpotence hypothesis of
\ref{lem:continuous-bar-cobar}. By
\ref{rmk:colim-does-not-recover-full} this is the
low-degree-sector failure; by
\ref{defn:nilpotent-low-degree-obstruction-terms} its obstruction
terms are
\(\mathfrak o^{\mathrm{lin}}_N\),
\(\mathfrak o^{\mathrm{ss}}\), and
\(\mathfrak o^{\mathrm{Cas}}_{\leq2}\); by
\ref{lem:tate-casimir-obstruction}
this is the same coefficient-category problem as
\ref{prob:formal-factorization-center}. Thus the nilpotent-truncation route
closes the classical nilpotent cotangent shadow on $\mathfrak m^3$, and the
residual moves from \ref{prob:formal-factorization-center} as stated
to its restricted form on the low-degree sector
$\mathfrak m/\mathfrak m^3=H_1\oplus H_2$ alone.
\end{prob}

\begin{thm}[Semisimple obstruction to lower-central Tate-pronilpotence]
\label{rmk:conilpotency-essential}
A second obstruction to extending $\Phi_{\mathfrak m^3}$ to all of
$\mathfrak h$ is structural rather than analytic: the inclusion of
$H_2=\mathfrak{sl}_2$ violates the lower-central Tate-pronilpotence route
used in Lemma~\ref{lem:continuous-bar-cobar} because
$[H_2,H_2]=H_2$ is its own descending Lie series. Even if the
degree-$1$ and degree-$2$ Casimir components converged in some
weighted Tate completion, the pronilpotent bar-cobar route would not
apply directly to the resulting Lie algebra.  This does not say that
ordinary finite-dimensional bar-cobar or PBW formalism fails for
\(\mathfrak{sl}_2\); it says that the manuscript's chosen lower-central
Tate-pronilpotent criterion must be replaced by a relative or reductive
model that builds in semisimple subalgebras separately.
\end{thm}

\begin{proof}
  Let \(L\) be any filtered Lie algebra containing the degree-two
  Hamiltonians \(H_2\) as a Lie subalgebra with the usual Poisson
  bracket.  The descending Lie powers of \(L\) satisfy
  \(L^{(n)}\cap H_2\supset H_2\) for every \(n\), because
  \([H_2,H_2]=H_2\).  Hence
  \[
    H_2\subset \bigcap_{n\geq1} L^{(n)} .
  \]
  The intersection of the descending Lie powers is therefore nonzero.
  This contradicts the lower-central Tate-pronilpotence condition used in
  Lemma~\ref{lem:continuous-bar-cobar}.  The bar-cobar argument may be
  applied to the nilpotent radical \(\mathfrak m^3\), but it cannot be
  applied unchanged to a Lie source containing the semisimple
  \(H_2\)-sector.
\end{proof}

\subsection*{Costello--Li precedent}

\begin{rmk}[Costello--Li balanced matrix truncation as structural
model]\label{rmk:costello-li-precedent}
Costello and Li \cite{costello-li-quantum-bcov} prove the
quantum perturbative BCOV theorem on flat Calabi--Yau by an
$N\to\infty$ stable connected limit through the open gauge algebra
$\mathfrak{gl}(N\mid N)$, with the directed Lie embedding
$\mathfrak{gl}(N\mid N)\hookrightarrow\mathfrak{gl}(N+k\mid N+k)$
and the BV-cohomology vanishing
$H^*_{\mathrm{BV}}\bigl(\mathfrak{gl}(N\mid N)\bigr)=0$
(supertrace identity $\mathrm{str}(I)=N-N=0$ kills the would-be
U(1)-anomaly). The structural analogy with the truncation here:
\begin{enumerate}
  \item Each window $\mathfrak{gl}(N\mid N)$ is finite-dimensional,
    BV-cohomology vanishes, the universal CE/Schouten identification
    and the BV-propagator coefficient kernel are well-defined on
    each window.
  \item The inverse / direct system is the Lie embedding
    $\mathfrak{gl}(N\mid N)\hookrightarrow\mathfrak{gl}(N+1\mid N+1)$
    on the Costello--Li side, the Lie surjection
    $\mathfrak h_{\leq N+1}\twoheadrightarrow\mathfrak h_{\leq N}$
    here.
  \item The U(1)-anomaly is killed by the structure of each window
    on the Costello--Li side (supertrace), and is removed by the
    $\mathfrak m^3$-truncation here (the cocycle $\omega$
    pairs only multidegree-$(1,0)$ with multidegree-$(0,1)$, both
    in $H_1\subset\mathfrak m\setminus\mathfrak m^2$, hence outside
    $\mathfrak m^3$).
  \item The colimit $N\to\infty$ converges on the Costello--Li side
    to the full $\mathfrak{gl}(\infty\mid\infty)$, providing the
    stable connected single-trace algebra. Here the inverse limit
    converges to $\mathfrak m^3\subsetneq\mathfrak h$, not to all
    of $\mathfrak h$, by \ref{rmk:colim-does-not-recover-full}.
\end{enumerate}
The structural difference: the
$\mathfrak{gl}(N\mid N)$ system is graded by matrix size, with
trivial action of one window on the next-window orthogonal
complement; the
$\mathfrak h=\mathfrak m$ system is graded by $\mathfrak m$-adic
degree, with the linear and quadratic sectors acting non-trivially
on every higher-degree window, which is precisely the
finite-window obstruction of
\ref{lem:finite-window-hamiltonian-obstruction}(iii) plus the
$\mathfrak{sl}_2\subset H_2$ semisimple obstruction to
lower-central Tate-pronilpotence.
\end{rmk}

\subsection*{Summary}

The nilpotent truncation
$\mathfrak h_{\leq N}=\mathfrak m^3/\mathfrak m^{N+1}$ closes
the classical finite nilpotent cotangent CE/PV shadow on the truncated
sector (degree-zero Hamiltonian plus finite cotangent labels). It does
not construct the unreduced open BV factorization-centre morphism. The quantum
all-order version on the same finite source is obstructed by
\ref{thm:phi-hbar-trunc}: the Moyal bracket leaves
$\mathfrak m^3/\mathfrak m^{N+1}$ for $N\geq4$.  The family
$\Phi_{\mathfrak m^3}=\{\Phi_{\leq N}\}_N$ of
\ref{prop:colim-recovery} is the classical coordinate-window cotangent
lift datum on the full classical $\mathfrak m^3$ sector; a completed
cochain inverse-limit statement requires the variance and
continuous-dualization datum named there. What is lost: the low-degree sector
$\mathfrak m/\mathfrak m^3=H_1\oplus H_2$ (linear translations and
the conformal-symplectic $\mathfrak{sl}_2$), the
constant-Hamiltonian $\hbar$-normalization, the U(1) center anomaly
$[\bar c]$, and the stable connected limit on the full
$\mathfrak h$. Recovering the low-degree sector requires the
coefficient-category change of
\ref{prob:formal-factorization-center} together with a
relative or reductive replacement for the lower-central
Tate-pronilpotent route of \ref{lem:continuous-bar-cobar}, so that the
semisimple $\mathfrak{sl}_2\subset H_2$ is treated as semisimple data
rather than as a pronilpotent radical. Thus the
nilpotent truncation
closes only the classical $\mathfrak m^3$ truncation; its all-order
quantum version is the separate Moyal-flat truncation problem above.

This is the same structural truncation as Costello--Li
$\mathfrak{gl}(N\mid N)\to\mathfrak{gl}(\infty\mid\infty)$, with the
key difference that the symplectic-plane $\mathfrak m$-adic system
does not admit a free directed exhaustion by nilpotent Lie
subalgebras containing $H_1\oplus H_2$; the truncation removes the
low-degree sector by hypothesis.
